% ============================================================
%  Resultados — Objetivo Específico 1 (NutriAPP)
%  Compilar con: pdflatex -> bibtex -> pdflatex -> pdflatex
%  Las figuras faltantes se muestran como recuadro gris con el
%  nombre de archivo esperado (ver imagenes/LEEME.md).
% ============================================================
\documentclass[12pt,a4paper]{article}

\usepackage[utf8]{inputenc}
\usepackage[T1]{fontenc}
\usepackage[spanish,es-tabla]{babel}
\usepackage[top=2.5cm,bottom=2.5cm,left=3cm,right=2.5cm]{geometry}

\usepackage{booktabs}
\usepackage{longtable}
\usepackage{array}
\usepackage{tabularx}
\usepackage{makecell}
\usepackage{graphicx}
\usepackage{float}
\usepackage{xcolor}
\usepackage{listings}
\usepackage{setspace}
\onehalfspacing
\usepackage{parskip}
\usepackage{enumitem}
\usepackage{caption}
\captionsetup{font=small, labelfont=bf}
\usepackage[hidelinks]{hyperref}
\usepackage{url}

\usepackage[backend=bibtex, style=ieee, sorting=none]{biblatex}
\addbibresource{../referencias.bib}
\addbibresource{referencias_obj1.bib}

% ── Numeración tipo memoria UNL: 6. / 6.1 / 6.1.1 / 6.1.1.1 ──
\setcounter{secnumdepth}{4}
\setcounter{tocdepth}{4}
\setcounter{section}{5}

% ── Rutas de imágenes ───────────────────────────────────────
\graphicspath{{imagenes/}{../documentacion_tecnica/imagenes/}}

% ── Figura con marcador si el archivo aún no existe ─────────
% \figura[ancho]{archivo}{pie de figura}{etiqueta}
\newcommand{\figura}[4][0.85\linewidth]{%
  \begin{figure}[H]
    \centering
    \IfFileExists{imagenes/#2}{\includegraphics[width=#1]{#2}}{%
      \IfFileExists{../documentacion_tecnica/imagenes/#2}{\includegraphics[width=#1]{#2}}{%
        \fbox{\parbox[c][4cm][c]{0.8\linewidth}{\centering\color{gray}%
          \textbf{FALTA IMAGEN}\\\texttt{imagenes/\detokenize{#2}}}}}}
    \caption{#3}
    \label{#4}
  \end{figure}}

% Marca visible para lo que aún debes completar
\newcommand{\pendiente}[1]{\textcolor{red}{\textbf{[PENDIENTE: #1]}}}

% ── Código TypeScript ───────────────────────────────────────
\definecolor{codebg}{RGB}{248,248,248}
\lstdefinelanguage{TypeScript}{
  keywords={const,let,function,return,export,import,from,async,await,
            try,catch,if,else,new,true,false,null,undefined},
  sensitive=true, morecomment=[l]{//}, morecomment=[s]{/*}{*/},
  morestring=[b]', morestring=[b]", morestring=[b]`,
}
\lstset{
  language=TypeScript, backgroundcolor=\color{codebg},
  basicstyle=\ttfamily\footnotesize, keywordstyle=\color{blue!70!black}\bfseries,
  stringstyle=\color{green!40!black}, commentstyle=\color{gray}\itshape,
  numbers=left, numberstyle=\tiny\color{gray}, frame=single,
  rulecolor=\color{gray!40}, breaklines=true, captionpos=b, tabsize=2,
  showstringspaces=false, extendedchars=true,
  literate={á}{{\'a}}1 {é}{{\'e}}1 {í}{{\'\i}}1 {ó}{{\'o}}1 {ú}{{\'u}}1
           {ñ}{{\~n}}1 {Á}{{\'A}}1 {É}{{\'E}}1 {Ó}{{\'O}}1 {×}{{$\times$}}1
           {—}{{---}}1,
}
\renewcommand{\lstlistingname}{Código}

\begin{document}

\section{Resultados}

% ============================================================
\subsection{Objetivo 1. Diseñar un prototipo funcional de sistema de gestión nutricional que integre datos antropométricos y bioquímicos del paciente con un catálogo de alimentos ecuatorianos, empleando un modelo de lenguaje largo con base de conocimiento nutricional estructurada}
% ============================================================

Para el desarrollo del Objetivo~1 se siguieron las fases de la metodología
RAD~\cite{fauzi2023} definidas en la sección de Metodología, organizadas en
cuatro actividades según el cronograma del proyecto:

\begin{enumerate}[leftmargin=2em]
  \item \textbf{Planificación de requisitos e identificación del catálogo
        alimentario ecuatoriano} (semanas 1--3): especificación de requisitos
        funcionales y no funcionales, identificación de actores, modelo
        conceptual y selección de las fuentes del catálogo de alimentos.
  \item \textbf{Definición de la arquitectura C4 y configuración del RAG}
        (semanas 3--7): diagramas de casos de uso, modelo de dominio y de
        comportamiento, arquitectura en los cuatro niveles del modelo
        C4~\cite{brown2024c4} y configuración del pipeline de generación
        aumentada por recuperación~\cite{lewis2020}.
  \item \textbf{Desarrollo de la interfaz del sistema} (semanas 6--9):
        construcción del backend, del frontend y de la base de datos, junto con
        las pruebas de caja negra de la API.
  \item \textbf{Documentación de resultados} (semanas 9--10): consolidación de
        los artefactos producidos, documentación técnica y manual de usuario.
\end{enumerate}

Como resultado se obtuvo el prototipo web \textbf{NutriAPP}, cuyos artefactos
se describen en las secciones siguientes.

% ============================================================
\subsubsection{Planificación de requisitos}
% ============================================================

\paragraph{Requisitos funcionales y no funcionales}

Los requisitos se obtuvieron a partir del levantamiento de información con el
actor principal del sistema, el nutricionista.
\pendiente{indica cómo se levantaron: reunión/entrevista con quién, fecha y si
fueron validados por el director o un nutricionista}.
La Tabla~\ref{tab:rf} presenta los doce requisitos funcionales y la
Tabla~\ref{tab:rnf} los siete requisitos no funcionales.

{\small
\begin{longtable}{|p{1.3cm}|p{12.2cm}|}
\caption{Requisitos funcionales del prototipo NutriAPP}\label{tab:rf}\\
\hline
\textbf{Código} & \textbf{Descripción} \\ \hline
\endfirsthead
\hline
\textbf{Código} & \textbf{Descripción} \\ \hline
\endhead
RF01 & El sistema debe permitir al nutricionista registrar un nuevo paciente con sus datos personales básicos (nombre, apellido, correo electrónico, teléfono, fecha de nacimiento y sexo). \\ \hline
RF02 & El sistema debe permitir al nutricionista registrar y actualizar el historial clínico del paciente, incluyendo antecedentes médicos, quirúrgicos y familiares, motivo de consulta, medicamentos, alergias e intolerancias alimentarias. \\ \hline
RF03 & El sistema debe permitir al nutricionista registrar medidas antropométricas del paciente: peso, talla, circunferencia de cintura y cadera, porcentaje de grasa corporal y masa muscular. \\ \hline
RF04 & El sistema debe calcular automáticamente el Índice de Masa Corporal (IMC) y la relación cintura-cadera, mostrando los resultados al nutricionista para su análisis clínico. \\ \hline
RF05 & El sistema debe permitir al nutricionista registrar datos biométricos del paciente: glucosa, HbA1c, colesterol total, triglicéridos, HDL, LDL, presión arterial y hemoglobina. \\ \hline
RF06 & El sistema debe permitir al nutricionista registrar los hábitos dietéticos del paciente, incluyendo frecuencia de comidas, consumo de agua, preferencias y restricciones alimentarias. \\ \hline
RF07 & El sistema debe calcular y proporcionar al nutricionista el Gasto Energético Total (GET) del paciente mediante la fórmula de Harris-Benedict, considerando sexo, peso, talla, edad y factor de actividad física seleccionado. \\ \hline
RF08 & El sistema debe calcular y proporcionar al nutricionista la distribución de macronutrientes diarios (proteínas, carbohidratos y grasas) en gramos y porcentajes, visualizando los resultados mediante un gráfico de dona interactivo. \\ \hline
RF09 & El sistema debe permitir al nutricionista interactuar con un asistente de inteligencia artificial mediante un chat individual por paciente, utilizando los datos clínicos, antropométricos y biométricos del paciente seleccionado como contexto. \\ \hline
RF10 & El sistema debe permitir al nutricionista consultar un chat global con el asistente de inteligencia artificial, el cual tiene acceso al contexto clínico de todos los pacientes registrados en el sistema. \\ \hline
RF11 & El sistema debe permitir al nutricionista cargar documentos PDF como base de conocimiento del asistente de inteligencia artificial, mediante un pipeline de generación aumentada por recuperación (RAG). \\ \hline
RF12 & El sistema debe permitir al nutricionista gestionar el catálogo de alimentos ecuatorianos almacenado en la base de datos, incluyendo nombre, grupo alimentario y composición nutricional por cada 100 gramos. \\ \hline
\end{longtable}
}

{\small
\begin{longtable}{|p{1.5cm}|p{2.6cm}|p{9.4cm}|}
\caption{Requisitos no funcionales del prototipo NutriAPP}\label{tab:rnf}\\
\hline
\textbf{Código} & \textbf{Categoría} & \textbf{Descripción} \\ \hline
\endfirsthead
\hline
\textbf{Código} & \textbf{Categoría} & \textbf{Descripción} \\ \hline
\endhead
RNF01 & Rendimiento & El sistema debe procesar y responder las consultas al asistente de inteligencia artificial en un tiempo máximo de 10 segundos bajo condiciones normales de conectividad. \\ \hline
RNF02 & Seguridad & El sistema debe garantizar la protección de las credenciales del nutricionista mediante cifrado de contraseñas con bcrypt y gestión de sesiones mediante tokens JWT en cabeceras HTTP. \\ \hline
RNF03 & Precisión & El margen de error en los cálculos automáticos de indicadores corporales (IMC, relación cintura-cadera) y gasto energético (Harris-Benedict) no debe exceder el 1\,\% respecto a los valores de referencia establecidos por las fórmulas oficiales. \\ \hline
RNF04 & Usabilidad & El sistema debe ofrecer una interfaz web clara e intuitiva, accesible desde navegadores modernos (Chrome, Firefox, Edge), sin necesidad de instalación adicional por parte del nutricionista. \\ \hline
RNF05 & Portabilidad & El sistema debe poder ejecutarse en un entorno de desarrollo local sin dependencia de infraestructura en la nube, utilizando únicamente la API de Groq para la inferencia del modelo de lenguaje durante la fase de prototipado. \\ \hline
RNF06 & Privacidad & Los datos clínicos procesados por el sistema corresponden a perfiles de prueba construidos con fines experimentales; en una implementación productiva, el sistema deberá cumplir con la Ley Orgánica de Protección de Datos Personales, garantizando el consentimiento explícito y el cifrado de la información del paciente. \\ \hline
RNF07 & Mantenibilidad & El código fuente del sistema debe estar versionado en GitHub con ramas \texttt{main} y \texttt{develop}, garantizando la trazabilidad del desarrollo y facilitando futuras actualizaciones. \\ \hline
\end{longtable}
}

\paragraph{Actores}

Se identificó un único actor humano, descrito en la Tabla~\ref{tab:actores}.
La API de Groq interviene como sistema externo de inferencia, no como actor
que inicie casos de uso.

\begin{table}[H]
\centering\small
\caption{Actores identificados}\label{tab:actores}
\begin{tabularx}{\linewidth}{|p{2.6cm}|X|}
\hline
\textbf{Actor} & \textbf{Descripción} \\ \hline
Nutricionista & Profesional de la salud que utiliza el sistema para registrar y gestionar los datos clínicos, antropométricos y biométricos de sus pacientes, consultar el asistente de inteligencia artificial y visualizar recomendaciones de planes nutricionales personalizados. \\ \hline
\end{tabularx}
\end{table}

\paragraph{Modelo conceptual}

A partir de los requisitos funcionales se elaboró el modelo conceptual de la
Figura~\ref{fig:conceptual}, que identifica las entidades principales del
dominio y sus relaciones. La Tabla~\ref{tab:glosario} define cada término.

\figura{modelo_conceptual.png}{Modelo conceptual de NutriAPP}{fig:conceptual}

\begin{table}[H]
\centering\small
\caption{Glosario de términos del modelo conceptual}\label{tab:glosario}
\begin{tabularx}{\linewidth}{|p{3.8cm}|X|}
\hline
\textbf{Término} & \textbf{Definición} \\ \hline
Nutricionista & Usuario autenticado que registra pacientes, valida planes y conversa con el asistente de IA. \\ \hline
Paciente & Persona atendida por el nutricionista; no posee credenciales de acceso al sistema. \\ \hline
HistorialClinico & Antecedentes médicos, quirúrgicos y familiares, motivo de consulta, medicación, alergias e intolerancias. \\ \hline
MedidasAntropometricas & Registro fechado de peso, talla, circunferencias, \% de grasa y masa muscular. \\ \hline
DatosBiometricos & Registro fechado de valores bioquímicos (glucosa, HbA1c, perfil lipídico, hemoglobina) y presión arterial. \\ \hline
HabitosDieteticos & Frecuencia de comidas, consumo de agua, preferencias y restricciones alimentarias. \\ \hline
SessionChat & Conversación del nutricionista con el asistente, individual por paciente o global. \\ \hline
MensajeChat & Mensaje del nutricionista o respuesta del asistente dentro de una sesión. \\ \hline
DocumentoConocimiento & Archivo PDF cargado por el nutricionista como contexto adicional del asistente. \\ \hline
PlanNutricional & Recomendación generada por el asistente y validada por el nutricionista. \\ \hline
Alimento & Elemento del catálogo con su grupo de equivalentes y composición nutricional. \\ \hline
\end{tabularx}
\end{table}

% ============================================================
\subsubsection{Identificación del catálogo de alimentos ecuatoriano}
% ============================================================

\paragraph{Fuentes de referencia}

Para construir un catálogo de alimentos adaptado al contexto ecuatoriano se
emplearon tres fuentes con funciones distintas (Tabla~\ref{tab:fuentes}).

\begin{table}[H]
\centering\small
\caption{Fuentes utilizadas para el catálogo de alimentos}\label{tab:fuentes}
\begin{tabularx}{\linewidth}{|p{4.2cm}|X|}
\hline
\textbf{Fuente} & \textbf{Uso en el catálogo} \\ \hline
ENSANUT~\cite{msp_ensanut} & Contexto epidemiológico de la malnutrición en Ecuador e identificación de alimentos representativos del consumo nacional. \\ \hline
Tabla de Composición de Alimentos, Cuenca~\cite{ortiz2018tabla} & Valores nutricionales por cada 100\,g de alimentos de producción local y preparaciones de la región andina sur (129 alimentos analizados entre 2011 y 2017). \\ \hline
SMAE, 4.ª ed.~\cite{perez2014smae} & Criterio de agrupación por equivalentes y porción sugerida, que permite intercambiar alimentos de un mismo grupo sin alterar de forma importante la energía y los macronutrientes del plan. \\ \hline
\end{tabularx}
\end{table}

La Tabla de Cuenca aporta la composición de los alimentos locales y el SMAE
aporta la estructura de grupos y porciones sobre la que el asistente arma las
recomendaciones (Figuras~\ref{fig:tabla_cuenca} y~\ref{fig:smae_frutas}).

\figura{tabla_cuenca_cereales.png}{Fragmento de la Tabla de Composición de Alimentos de Cuenca, grupo Cereales y derivados~\cite{ortiz2018tabla}}{fig:tabla_cuenca}
\figura{smae_frutas.png}{Fragmento del grupo Frutas del SMAE~\cite{perez2014smae}}{fig:smae_frutas}

\paragraph{Procesamiento e incorporación al sistema}

El contenido del SMAE se extrajo del documento fuente mediante scripts de
conversión y se generó un archivo de carga inicial
(\texttt{backend/prisma/smae\_seed\_data.ts}) con 2392 alimentos. Cada registro
se almacena en la tabla \texttt{foods} con los campos de la
Tabla~\ref{tab:campos_food}; el campo \texttt{source} distingue los alimentos
cargados del SMAE de los registrados manualmente por el nutricionista (RF12).

\pendiente{describe cómo se adaptaron los alimentos locales: cuántos alimentos
de la Tabla de Cuenca se añadieron o sustituyeron, con qué criterio, y si se
marcaron con otro valor en \texttt{source}. Hoy la carga inicial sólo contiene
registros SMAE}.

\begin{table}[H]
\centering\small
\caption{Estructura de un alimento en el catálogo}\label{tab:campos_food}
\begin{tabularx}{\linewidth}{|p{3.6cm}|X|}
\hline
\textbf{Campo} & \textbf{Descripción} \\ \hline
\texttt{name}, \texttt{description} & Nombre y descripción del alimento. \\ \hline
\texttt{grupoSmae}, \texttt{subgrupoSmae} & Grupo y subgrupo de equivalentes (p.\,ej. AOA / muy bajo aporte de grasa). \\ \hline
\texttt{porcionSugerida}, \texttt{unidadPorcion} & Cantidad de un equivalente y su unidad (taza, pieza, g). \\ \hline
\texttt{grossWeight}, \texttt{netWeight} & Peso bruto y neto de la porción (g). \\ \hline
\texttt{energyKcal}, \texttt{energyKj} & Energía de la porción. \\ \hline
\texttt{protein}, \texttt{fats}, \texttt{carbohydrates}, \texttt{fiber} & Macronutrientes y fibra (g). \\ \hline
\texttt{indiceGlucemico}, \texttt{cargaGlucemica} & Índice y carga glucémica, cuando la fuente los reporta. \\ \hline
\texttt{source} & Origen del registro (\texttt{SMAE} o manual). \\ \hline
\end{tabularx}
\end{table}

\paragraph{Distribución del catálogo}

La Tabla~\ref{tab:grupos} resume el número de alimentos por grupo de
equivalentes.

\begin{table}[H]
\centering\small
\caption{Alimentos del catálogo por grupo de equivalentes}\label{tab:grupos}
\begin{tabular}{|l|r|}
\hline
\textbf{Grupo} & \textbf{N.º de alimentos} \\ \hline
Verduras & \pendiente{n} \\ \hline
Frutas & \pendiente{n} \\ \hline
Cereales y tubérculos & \pendiente{n} \\ \hline
Leguminosas & \pendiente{n} \\ \hline
Alimentos de origen animal (AOA) & \pendiente{n} \\ \hline
Leche & \pendiente{n} \\ \hline
Aceites y grasas & \pendiente{n} \\ \hline
Azúcares & \pendiente{n} \\ \hline
Alimentos libres en energía & \pendiente{n} \\ \hline
\textbf{Total} & \textbf{\pendiente{n}} \\ \hline
\end{tabular}
\end{table}

% ============================================================
\subsubsection{Diseño del sistema}
% ============================================================

\paragraph{Modelo de casos de uso}

La Figura~\ref{fig:cu_general} presenta el diagrama general de casos de uso,
derivado de los requisitos funcionales. La Tabla~\ref{tab:traza_cu} muestra la
correspondencia entre casos de uso y requisitos.

\figura{casos_uso_general.png}{Diagrama general de casos de uso}{fig:cu_general}

\begin{table}[H]
\centering\small
\caption{Trazabilidad entre casos de uso y requisitos funcionales}\label{tab:traza_cu}
\begin{tabular}{|l|l|l|}
\hline
\textbf{Código} & \textbf{Caso de uso} & \textbf{RF} \\ \hline
CU01 & Iniciar sesión & RNF02 \\ \hline
CU02 & Gestionar paciente & RF01 \\ \hline
CU03 & Registrar historial clínico & RF02 \\ \hline
CU04 & Registrar medidas antropométricas & RF03, RF04 \\ \hline
CU05 & Registrar datos biométricos & RF05 \\ \hline
CU06 & Registrar hábitos dietéticos & RF06 \\ \hline
CU07 & Calcular GET y distribución de macronutrientes & RF07, RF08 \\ \hline
CU08 & Consultar asistente IA por paciente & RF09 \\ \hline
CU09 & Consultar asistente IA global & RF10 \\ \hline
CU10 & Cargar documento PDF & RF11 \\ \hline
CU11 & Gestionar catálogo de alimentos & RF12 \\ \hline
\end{tabular}
\end{table}
\pendiente{ajusta los códigos y nombres a los de tu diagrama}

Se describen a continuación los casos de uso centrales del prototipo, con su
flujo normal y sus flujos alternos. \pendiente{revisa que los mensajes coincidan
con los que muestra tu interfaz}

% ── Plantilla de ficha de caso de uso ───────────────────────
\newcommand{\fichaCU}[6]{%
\begin{longtable}{|p{3cm}|p{11cm}|}
\caption{Descripción del caso de uso #1}\\ \hline
\textbf{Caso de uso} & #1 \\ \hline
\textbf{Código} & #2 \\ \hline
\textbf{Actor} & Nutricionista \\ \hline
\textbf{Descripción} & #3 \\ \hline
\textbf{Precondiciones} & #4 \\ \hline
\multicolumn{2}{|c|}{\textbf{CURSO NORMAL DE EVENTOS}} \\ \hline
\multicolumn{2}{|p{14.5cm}|}{#5} \\ \hline
\multicolumn{2}{|c|}{\textbf{CURSO ALTERNO DE EVENTOS}} \\ \hline
\multicolumn{2}{|p{14.5cm}|}{#6} \\ \hline
\end{longtable}}

{\small
\fichaCU{Registrar medidas antropométricas}{CU04}
{El nutricionista registra las medidas antropométricas de un paciente y el sistema calcula el IMC y la relación cintura-cadera.}
{Sesión iniciada. Paciente seleccionado.}
{\begin{enumerate}[nosep]
 \item El nutricionista abre la sección de medidas antropométricas del paciente.
 \item Ingresa peso, talla, circunferencias de cintura y cadera, \% de grasa y masa muscular.
 \item El sistema calcula el IMC y la relación cintura-cadera y muestra su clasificación.
 \item El nutricionista guarda el registro.
 \item El sistema almacena la medición con su fecha y actualiza el historial del paciente.
\end{enumerate}}
{\textbf{A. Campos obligatorios vacíos o no numéricos:} el sistema muestra un aviso por campo y no guarda.\\
 \textbf{B. Error del servidor:} el sistema muestra un mensaje de error y conserva los datos del formulario.}

\fichaCU{Consultar asistente IA por paciente}{CU08}
{El nutricionista conversa con el asistente usando como contexto los datos clínicos del paciente seleccionado y el catálogo SMAE.}
{Sesión iniciada. Paciente seleccionado.}
{\begin{enumerate}[nosep]
 \item El nutricionista escribe una consulta en el chat del paciente.
 \item El sistema arma el contexto con los datos clínicos, antropométricos y biométricos del paciente.
 \item Si la consulta trata sobre planes, alimentos o equivalentes, el sistema añade al contexto una muestra del catálogo SMAE por grupo.
 \item El sistema envía al modelo de lenguaje el prompt del sistema, el contexto, el historial reciente y la consulta.
 \item El sistema muestra la respuesta en formato de tablas y guarda ambos mensajes en el historial.
\end{enumerate}}
{\textbf{A. Mensaje vacío:} el sistema responde «El mensaje es requerido».\\
 \textbf{B. Consulta fuera del dominio nutricional:} el asistente responde con el mensaje de rechazo definido en el prompt del sistema.\\
 \textbf{C. Límite de peticiones de la API:} el sistema informa que el servicio está saturado y pide reintentar.}

\fichaCU{Cargar documento PDF}{CU10}
{El nutricionista adjunta un documento PDF para que el asistente responda con base en su contenido.}
{Sesión iniciada.}
{\begin{enumerate}[nosep]
 \item El nutricionista adjunta un archivo PDF o de texto y escribe una consulta.
 \item El sistema valida el tipo y el tamaño del archivo (máximo 10\,MB).
 \item El sistema extrae el texto del documento.
 \item El sistema envía al modelo el texto extraído junto con la consulta.
 \item El sistema muestra la respuesta basada en el documento.
\end{enumerate}}
{\textbf{A. Tipo de archivo no permitido:} el sistema rechaza la carga.\\
 \textbf{B. Archivo mayor a 10\,MB:} el sistema rechaza la carga.\\
 \textbf{C. PDF sin texto extraíble (escaneado):} \pendiente{describe el comportamiento observado}.}
}

\paragraph{Modelo de dominio y comportamiento}

El modelo conceptual se refinó en el diagrama de clases de la
Figura~\ref{fig:clases}, con atributos, tipos y cardinalidades. Los diagramas
de secuencia de las Figuras~\ref{fig:seq_chat} y~\ref{fig:seq_pdf} describen los
dos flujos principales del asistente.

\figura{diagrama_clases.png}{Diagrama de clases de NutriAPP}{fig:clases}
\figura{secuencia_chat_paciente.png}{Diagrama de secuencia «Consultar asistente IA por paciente»}{fig:seq_chat}
\figura{secuencia_carga_pdf.png}{Diagrama de secuencia «Cargar documento PDF»}{fig:seq_pdf}

% ============================================================
\subsubsection{Arquitectura C4 y configuración del RAG}
% ============================================================

La arquitectura se documentó con el modelo C4~\cite{brown2024c4}, que describe
el sistema en cuatro niveles de abstracción.

\paragraph{Nivel 1: Contexto}

NutriAPP es el sistema central (Figura~\ref{fig:c4_contexto}). El nutricionista
accede desde un navegador web y la única dependencia externa es la API de Groq,
que provee la inferencia del modelo de lenguaje (RNF05).

\figura{c4_contexto.png}{Diagrama C4, nivel 1: contexto del sistema}{fig:c4_contexto}

\paragraph{Nivel 2: Contenedores}

El sistema se compone de tres contenedores (Figura~\ref{fig:c4_contenedores}):
una aplicación de página única en React servida por Vite, una API REST en
Express con TypeScript y una base de datos PostgreSQL accedida mediante Prisma ORM.

\figura{c4_contenedores.png}{Diagrama C4, nivel 2: contenedores}{fig:c4_contenedores}

\paragraph{Nivel 3: Componentes del backend}

El backend se organiza en cuatro routers con su controlador
(Tabla~\ref{tab:c4_comp} y Figura~\ref{fig:c4_componentes}). El servicio
\texttt{aiService} concentra la comunicación con la API de Groq y el middleware
de autenticación verifica el token JWT en cada ruta protegida (RNF02).

\begin{table}[H]
\centering\small
\caption{Componentes del backend (nivel 3 C4)}\label{tab:c4_comp}
\begin{tabularx}{\linewidth}{|l|l|X|}
\hline
\textbf{Router} & \textbf{Controlador} & \textbf{Responsabilidad} \\ \hline
\texttt{/api/auth} & \texttt{authController} & Inicio de sesión y administración de nutricionistas \\ \hline
\texttt{/api/patients} & \texttt{patientController} & Pacientes, historia clínica, biometría, antropometría, hábitos dietéticos y menú semanal \\ \hline
\texttt{/api/foods} & \texttt{foodController} & Catálogo de alimentos y grupos SMAE \\ \hline
\texttt{/api/ai} & \texttt{aiController} & Chat, recomendaciones, análisis de alimentos, planes de comida, carga de documentos e historial \\ \hline
\end{tabularx}
\end{table}

\figura{c4_componentes.png}{Diagrama C4, nivel 3: componentes del backend}{fig:c4_componentes}

\paragraph{Nivel 4: Código}

La Figura~\ref{fig:c4_codigo} detalla el servicio \texttt{aiService}, que expone
los métodos \texttt{chat}, \texttt{chatWithDocument},
\texttt{generateNutritionRecommendation}, \texttt{generateMealPlan} y
\texttt{analyzeFood}, además de la función interna \texttt{buildSmaeContext}.

\figura{c4_codigo.png}{Diagrama C4, nivel 4: código del servicio de IA}{fig:c4_codigo}

\paragraph{Configuración del pipeline RAG}

El asistente combina tres fuentes de contexto antes de cada llamada al modelo:

\begin{enumerate}[leftmargin=2em]
  \item \textbf{Datos clínicos del paciente}: historial, antropometría,
        biometría y hábitos del paciente seleccionado (RF09) o de todos los
        pacientes en el chat global (RF10).
  \item \textbf{Catálogo SMAE}: cuando la consulta contiene términos como
        \emph{plan}, \emph{alimento}, \emph{menú}, \emph{equivalente} o
        \emph{kcal}, el sistema recupera de la base de datos ocho alimentos por
        cada uno de siete grupos de equivalentes (56 alimentos) y los añade al
        prompt con su porción, energía, macronutrientes e índice glucémico.
  \item \textbf{Documento del nutricionista}: el texto extraído del PDF cargado
        (RF11), limitado a los primeros 6000 caracteres.
\end{enumerate}

\pendiente{precisa en el texto qué entiendes por «recuperación»: hoy el sistema
selecciona alimentos por grupo y trunca el PDF, no hace búsqueda semántica con
embeddings. Descríbelo tal cual o implementa la búsqueda antes de redactarlo}.

La Tabla~\ref{tab:rag_param} resume los parámetros configurados.

\begin{table}[H]
\centering\small
\caption{Parámetros de configuración del asistente de IA}\label{tab:rag_param}
\begin{tabularx}{\linewidth}{|p{5cm}|X|}
\hline
\textbf{Parámetro} & \textbf{Valor} \\ \hline
Proveedor de inferencia & API de Groq (compatible con OpenAI) \\ \hline
Modelo por defecto & \pendiente{confirma: el código usa \texttt{openai/gpt-oss-120b}; la metodología dice Llama~3.3~70B} \\ \hline
Modelos seleccionables & \texttt{gpt-oss-120b}, \texttt{llama-3.3-70b-versatile}, \texttt{llama-3.1-8b-instant} \\ \hline
Temperatura & 0.3 \\ \hline
Máximo de tokens de respuesta & 8000 \\ \hline
Historial enviado al modelo & 20 mensajes de la conversación \\ \hline
Contexto SMAE & 7 grupos $\times$ 8 alimentos, activado por palabras clave \\ \hline
Documentos admitidos & PDF o texto plano, hasta 10\,MB \\ \hline
Texto del documento usado & primeros 6000 caracteres \\ \hline
Extracción de texto & \texttt{pdf-parse} 2.x \\ \hline
\end{tabularx}
\end{table}

El comportamiento del asistente se restringe mediante un prompt del sistema
(\texttt{backend/prompts/system.md}) que define su alcance clínico, rechaza
solicitudes fuera del dominio nutricional, obliga a presentar los valores en
tablas y declara las fórmulas y referencias admitidas (Mifflin-St~Jeor,
Harris-Benedict, FAO/OMS, SMAE 4.ª ed.).

\begin{lstlisting}[caption={Construcción del contexto SMAE por grupo (extracto de \texttt{aiService.ts})},label={lst:smae}]
async function buildSmaeContext(grupos?: string[]): Promise<string> {
  const targetGroups = grupos ?? [
    'Verduras', 'Frutas', 'Cereales', 'Leguminosas', 'AOA', 'Leche', 'Aceites y grasas',
  ];
  const perGroup = 8;
  const perGroupResults = await Promise.all(
    targetGroups.map(g =>
      prisma.food.findMany({
        where: { source: 'SMAE', grupoSmae: g },
        select, orderBy: { name: 'asc' }, take: perGroup,
      })
    )
  );
  // ... se formatea cada alimento como:
  // "- Nombre (porción): kcal P:g G:g HC:g Fib:g IG"
}
\end{lstlisting}

% ============================================================
\subsubsection{Desarrollo de la interfaz del sistema}
% ============================================================

\paragraph{Tecnologías}

Las Tablas~\ref{tab:tec_back} y~\ref{tab:tec_front} presentan las tecnologías
empleadas y sus versiones, tomadas de los archivos \texttt{package.json} del
proyecto.

\begin{table}[H]
\centering\small
\caption{Tecnologías y librerías principales del backend}\label{tab:tec_back}
\begin{tabular}{|l|l|l|}
\hline
\textbf{Librería} & \textbf{Versión} & \textbf{Propósito} \\ \hline
Node.js & \pendiente{v} & Entorno de ejecución \\ \hline
TypeScript & \^{}5.4.5 & Tipado estático \\ \hline
Express & \^{}4.22.2 & API REST \\ \hline
Prisma (\texttt{@prisma/client}) & \^{}5.14.0 & ORM sobre PostgreSQL \\ \hline
PostgreSQL & \pendiente{v} & Base de datos relacional \\ \hline
openai (SDK, apuntando a Groq) & \^{}6.37.0 & Cliente de la API de inferencia \\ \hline
pdf-parse & \^{}2.4.5 & Extracción de texto de PDF \\ \hline
multer & \^{}2.2.0 & Recepción de archivos \\ \hline
jsonwebtoken & \^{}9.0.2 & Tokens de sesión JWT \\ \hline
bcryptjs & \^{}2.4.3 & Cifrado de contraseñas \\ \hline
helmet & \^{}8.2.0 & Cabeceras de seguridad HTTP \\ \hline
\end{tabular}
\end{table}

\begin{table}[H]
\centering\small
\caption{Tecnologías y librerías principales del frontend}\label{tab:tec_front}
\begin{tabular}{|l|l|l|}
\hline
\textbf{Librería} & \textbf{Versión} & \textbf{Propósito} \\ \hline
React & \^{}18.3.1 & Construcción de la interfaz \\ \hline
Vite & \^{}5.3.1 & Empaquetado y servidor de desarrollo \\ \hline
TailwindCSS & \^{}3.4.4 & Estilos \\ \hline
axios & \^{}1.7.2 & Consumo de la API REST \\ \hline
react-markdown + remark-gfm & \^{}10.1.0 / \^{}4.0.1 & Renderizado de respuestas del asistente con tablas \\ \hline
lucide-react & \^{}1.48.0 & Iconografía \\ \hline
\end{tabular}
\end{table}

\paragraph{Estructura del proyecto}

El backend sigue una organización por capas: \texttt{routes} define los
endpoints, \texttt{controllers} contiene la lógica de cada recurso,
\texttt{services} encapsula la integración con el modelo de lenguaje,
\texttt{middleware} gestiona autenticación y errores, y \texttt{prisma} contiene
el esquema, las migraciones y la carga inicial del catálogo
(Figura~\ref{fig:estructura_back}). El frontend organiza sus componentes por
módulo funcional en \texttt{components/features} (pacientes, historia clínica,
antropometría, biometría, hábitos dietéticos, alimentos, menús, progreso y
administración) (Figura~\ref{fig:estructura_front}).

\figura[0.4\linewidth]{estructura_backend.png}{Estructura del backend}{fig:estructura_back}
\figura[0.4\linewidth]{estructura_frontend.png}{Estructura del frontend}{fig:estructura_front}

\paragraph{Modelo de base de datos}

El esquema de Prisma define diez entidades: \texttt{User}, \texttt{Patient},
\texttt{ClinicalHistory}, \texttt{Biometrics}, \texttt{Anthropometry},
\texttt{DietaryHabits}, \texttt{ChatMessage}, \texttt{Food},
\texttt{WeeklyMenu} y \texttt{WeeklyMenuItem} (Figura~\ref{fig:er}).

\figura{diagrama_er.png}{Modelo entidad-relación de la base de datos}{fig:er}

\paragraph{Trazabilidad de requisitos en la implementación}

La Tabla~\ref{tab:traza_impl} relaciona cada requisito funcional con el
componente que lo implementa.

\begin{table}[H]
\centering\footnotesize
\caption{Requisitos funcionales y componente que los implementa}\label{tab:traza_impl}
\begin{tabularx}{\linewidth}{|l|X|X|}
\hline
\textbf{RF} & \textbf{Endpoint} & \textbf{Componente de interfaz} \\ \hline
RF01 & \texttt{POST/PUT /api/patients} & \texttt{features/patients} \\ \hline
RF02 & \texttt{PUT /api/patients/:id/clinical-history} & \texttt{features/clinical} \\ \hline
RF03 & \texttt{POST /api/patients/:id/anthropometry} & \texttt{features/anthropometry} \\ \hline
RF04 & --- (cálculo en cliente) & \texttt{utils/nutritionCalculations.ts} \\ \hline
RF05 & \texttt{POST /api/patients/:id/biometrics} & \texttt{features/biometrics} \\ \hline
RF06 & \texttt{POST /api/patients/:id/dietary-habits} & \texttt{features/dietary} \\ \hline
RF07 & --- (cálculo en cliente) & \texttt{NutritionistPanel.tsx} \\ \hline
RF08 & --- (cálculo en cliente) & \pendiente{componente del gráfico de dona} \\ \hline
RF09 & \texttt{POST /api/ai/chat} (con \texttt{patientId}) & \texttt{NutritionistPanel.tsx} \\ \hline
RF10 & \texttt{POST /api/ai/chat} (global) & \texttt{NutritionistPanel.tsx} \\ \hline
RF11 & \texttt{POST /api/ai/chat-upload} & \texttt{NutritionistPanel.tsx} \\ \hline
RF12 & \texttt{GET/POST/PUT/DELETE /api/foods} & \texttt{features/foods} \\ \hline
\end{tabularx}
\end{table}

\paragraph{Interfaz del sistema}

Las Figuras~\ref{fig:ui_login} a~\ref{fig:ui_pdf} muestran los módulos
principales del prototipo en funcionamiento.

\figura{im02_login.png}{Pantalla de inicio de sesión}{fig:ui_login}
\figura{im01_dashboard.png}{Panel principal del nutricionista}{fig:ui_dashboard}
\figura{im03_pacientes.png}{Listado de pacientes}{fig:ui_pacientes}
\figura{im04_registro_paciente.png}{Registro de paciente}{fig:ui_registro}
\figura{im05_antropometria.png}{Registro de medidas antropométricas con cálculo de IMC e ICC}{fig:ui_antropo}
\figura{im06_biometria.png}{Registro de datos biométricos}{fig:ui_bio}
\figura{im07_chat_ia.png}{Asistente de IA con contexto del paciente}{fig:ui_chat}
\figura{im08_chat_pdf.png}{Consulta al asistente con un documento PDF}{fig:ui_pdf}

\pendiente{añade capturas del cálculo de GET, del gráfico de macronutrientes
(RF07, RF08) y del catálogo de alimentos (RF12) si existen en la interfaz}

El código fuente se encuentra en el repositorio
\pendiente{URL del repositorio de GitHub}.

% ============================================================
\subsubsection{Pruebas de caja negra}
% ============================================================

Los endpoints de la API se validaron con pruebas de caja negra ejecutadas en
Thunder Client. Cada caso define la petición, el código HTTP esperado y, cuando
aplica, un texto que debe aparecer en la respuesta. La colección completa se
encuentra en el archivo \texttt{NutriApp\_ThunderClient.json}
\pendiente{referencia al anexo}. La Tabla~\ref{tab:pcn} resume los resultados.

\pendiente{ejecuta la colección y llena la columna «Estado» con lo que
realmente obtuviste; anota los casos fallidos y qué se corrigió}

{\footnotesize
\begin{longtable}{|p{1.2cm}|p{2.9cm}|p{5.6cm}|p{2.6cm}|p{1.6cm}|}
\caption{Pruebas de caja negra de la API REST}\label{tab:pcn}\\
\hline
\textbf{ID} & \textbf{Endpoint} & \textbf{Descripción} & \textbf{Esperado} & \textbf{Estado} \\ \hline
\endfirsthead
\hline
\textbf{ID} & \textbf{Endpoint} & \textbf{Descripción} & \textbf{Esperado} & \textbf{Estado} \\ \hline
\endhead
TC-01 & GET /health & Verificar que el servidor está disponible & 200, «ok» & \pendiente{} \\ \hline
TC-02 & POST /api/auth/login & Rechazar inicio de sesión sin cuerpo & 400, «requeridos» & \pendiente{} \\ \hline
TC-03 & POST /api/auth/login & Rechazar contraseña incorrecta & 401, «incorrectos» & \pendiente{} \\ \hline
TC-04 & GET /api/auth/me & Impedir acceso a ruta protegida sin token & 401, «Token requerido» & \pendiente{} \\ \hline
TC-05 & GET /api/auth/me & Rechazar token malformado & 401, «inválido» & \pendiente{} \\ \hline
TC-06 & GET /api/auth/me & Token válido de usuario inexistente en BD & 200 & \pendiente{} \\ \hline
TC-07 & GET /api/patients & Listar pacientes registrados & 200, arreglo & \pendiente{} \\ \hline
TC-08 & POST /api/patients & Rechazar paciente con campos faltantes & 400, «campos requeridos» & \pendiente{} \\ \hline
TC-09 & GET /api/patients/:id & Paciente con ID inexistente & 404, «no encontrado» & \pendiente{} \\ \hline
TC-10 & GET /api/foods & Listar alimentos del catálogo SMAE & 200, arreglo & \pendiente{} \\ \hline
TC-11 & GET /api/foods/smae-groups & Obtener grupos SMAE disponibles & 200 & \pendiente{} \\ \hline
TC-12 & GET /api/foods/:id & Alimento con ID inexistente & 404 & \pendiente{} \\ \hline
TC-13 & POST /api/ai/chat & Rechazar mensaje vacío & 400, «requerido» & \pendiente{} \\ \hline
TC-14 & POST /api/ai/chat & Rechazar petición sin campo \texttt{message} & 400 & \pendiente{} \\ \hline
TC-15 & POST /api/ai/recommendation & Rechazar recomendación con datos del paciente faltantes & 400 & \pendiente{} \\ \hline
TC-16 & POST /api/ai/analyze-food & Rechazar análisis sin campos requeridos & 400 & \pendiente{} \\ \hline
TC-17 & POST /api/ai/meal-plan & Rechazar plan sin calorías objetivo & 400 & \pendiente{} \\ \hline
TC-18 & GET /api/ai/chat-history/global & Obtener historial del chat global & 200, arreglo & \pendiente{} \\ \hline
TC-19 & POST /api/ai/chat & Solicitar cálculo de GET nutricional & 200, «kcal» & \pendiente{} \\ \hline
TC-20 & POST /api/ai/chat & Rechazar solicitud fuera de dominio (código Python) & 200, «Solo puedo asistirte» & \pendiente{} \\ \hline
TC-21 & POST /api/ai/recommendation & Recomendación completa con equivalentes SMAE & 200, «SMAE» & \pendiente{} \\ \hline
TC-22 & POST /api/ai/analyze-food & Análisis de un alimento SMAE (tortilla de maíz) & 200, «kcal» & \pendiente{} \\ \hline
\end{longtable}
}

\figura{im10_postman_rag.png}{Ejemplo de ejecución de una prueba sobre el asistente de IA}{fig:prueba_rag}

\paragraph{Verificación de requisitos no funcionales}

\pendiente{opcional pero recomendable: mide y reporta RNF01 (tiempo de respuesta
del asistente, p.\,ej. promedio de 10 consultas) y RNF03 (compara IMC, ICC y GET
del sistema contra cálculo manual en 3--5 pacientes de prueba)}

\begin{table}[H]
\centering\small
\caption{Verificación de requisitos no funcionales}\label{tab:rnf_verif}
\begin{tabularx}{\linewidth}{|l|X|X|}
\hline
\textbf{RNF} & \textbf{Criterio} & \textbf{Resultado} \\ \hline
RNF01 & Respuesta del asistente $\leq$ 10\,s & \pendiente{tiempo medido} \\ \hline
RNF02 & Contraseñas con bcrypt y sesión con JWT & Verificado en TC-02 a TC-05 \\ \hline
RNF03 & Error $\leq$ 1\,\% en IMC, ICC y GET & \pendiente{error medido} \\ \hline
RNF04 & Funciona en Chrome, Firefox y Edge & \pendiente{} \\ \hline
RNF05 & Ejecución local con Groq como único servicio externo & \pendiente{} \\ \hline
RNF07 & Repositorio con ramas \texttt{main} y \texttt{develop} & \pendiente{} \\ \hline
\end{tabularx}
\end{table}


\printbibliography[title={Referencias de esta sección}]

\end{document}
